\chapter{The Euler--Lagrange equation and the smooth bootstrap}
\label{ch:euler-lagrange}
\module{NoCompromise.Energy.PotentialHolder,
        NoCompromise.Stationary.Bootstrap,
        NoCompromise.Stationary.Connected,
        NoCompromise.Stationary.Defs,
        NoCompromise.Stationary.EulerLagrange,
        NoCompromise.Stationary.MinimizerContext,
        NoCompromise.Stationary.ScalingIdentity,
        NoCompromise.Variation.CoulombBoundary}

\begin{notation}\label{not:minimizer-rep}
\uses{def:minimizer,conv:representative,prop:no-singular-points}
For the rest of Part~\ref{part:regularity}, $E$ is a fixed-volume minimiser at
volume $V>0$ and $\Omega:=E^{(1)}$ is its bounded open representative with
compact $C^{1,1/2}$ boundary, as produced by
Lemmas~\ref{lem:quasiminimal}, \ref{lem:bounded-representative} and
Proposition~\ref{prop:no-singular-points}.  $\nu$ denotes the continuous
outward normal (Convention~\ref{conv:outward}).
\end{notation}

\section{The nonlocal first variation}

\begin{lemma}[First variation of the Coulomb energy]\label{lem:coulomb-first-var}
\uses{not:minimizer-rep,def:coulomb,def:straight}
Let $\Omega$ be a bounded $C^1$ domain and $X\in C^\infty_c(\R^3;\R^3)$.  Then
\begin{equation}\label{eq:coulomb-first-var}
  \frac{d}{dt}\Big|_{t=0}D\bigl(F^X_t(\Omega)\bigr)
  =\int_{\partial\Omega}v_\Omega\,X\cdot\nu\,d\Hs^2 .
\end{equation}
\end{lemma}

\begin{proof}
\uses{cor:W11-GG-uses,lem:straight-diffeo,def:coulomb,lem:straight-cofactor}
Pull both integrations back to $\Omega$:
\[
  D\bigl(F^X_t(\Omega)\bigr)
  =\frac12\iint_{\Omega\times\Omega}
   \frac{JF^X_t(x)\,JF^X_t(y)}{|F^X_t(x)-F^X_t(y)|}\,dx\,dy .
\]
Differentiating the Jacobians and the kernel gives an integrand bounded by
\[
  C\Bigl(\frac1{|x-y|}+\frac{|X(x)-X(y)|}{|x-y|^2}\Bigr)
  \le\frac{C'}{|x-y|},
\]
since $X$ is Lipschitz; this majorant is integrable on the bounded product
$\Omega\times\Omega$.  Dominated convergence and exchange of $x,y$ yield
\begin{equation}\label{eq:coulomb-first-var-steps}
\begin{aligned}
  \frac{d}{dt}\Big|_0D\bigl(F^X_t(\Omega)\bigr)
  &=\iint_{\Omega\times\Omega}
    \Bigl(\frac{\dvg X(x)}{|x-y|}-\frac{(x-y)\cdot X(x)}{|x-y|^3}\Bigr)dx\,dy\\
  &=\iint_{\Omega\times\Omega}
    \dvg_x\Bigl(\frac{X(x)}{|x-y|}\Bigr)dx\,dy .
\end{aligned}
\end{equation}
For the last step, fix $y\in\Omega$ outside a null set and apply
Corollary~\ref{cor:W11-GG-uses}\ref{it:W11-kernel} to
$W_y(x):=X(x)/|x-y|$, obtaining \eqref{eq:W11-GG-application}.  Both sides of
\eqref{eq:W11-GG-application} are integrable in $y$: the volume integrand is
bounded by $C(|x-y|^{-1}+|x-y|^{-2})$ and the boundary kernel is integrable on
$\partial\Omega\times\Omega$ in local graph coordinates.  Fubini gives
\eqref{eq:coulomb-first-var}.
\end{proof}

\begin{fnote}\label{fnote:C1-suffices}
The boundary is only $C^{1,1/2}$ at this stage, so the proof must not use a
smooth tubular normal graph.  Straight ambient perturbations
(Definition~\ref{def:straight}) suffice, and
Corollary~\ref{cor:W11-GG-uses} is stated for $C^1$ domains for exactly this
reason.
\end{fnote}

\begin{theorem}[Constrained first variation]\label{thm:constrained-first-var}
\uses{not:minimizer-rep,lem:coulomb-first-var,cor:W11-GG-uses}
There is $\lambda\in\R$ such that for every $X\in C^\infty_c(\R^3;\R^3)$,
\begin{equation}\label{eq:weak-EL}
  \int_{\partial\Omega}\dvgt X\,d\Hs^2
  =\int_{\partial\Omega}(\lambda-v_\Omega)\,X\cdot\nu\,d\Hs^2 .
\end{equation}
Explicitly, with
\[
  A(X):=\int_{\partial\Omega}\bigl(\dvgt X+v_\Omega X\cdot\nu\bigr),
  \qquad
  B(X):=\int_{\partial\Omega}X\cdot\nu ,
\]
one has $A(X)=\lambda B(X)$ with $\lambda=A(Y)/B(Y)$ for any $Y$ as below.
\end{theorem}

\begin{proof}
\uses{thm:first-var-perimeter,thm:first-var-volume,lem:coulomb-first-var,
      cor:W11-GG-uses,def:minimizer,lem:straight-diffeo}
Choose a compactly supported smooth field $Y$ equal to $x$ on a neighbourhood
of $\overline\Omega$.  By
Corollary~\ref{cor:W11-GG-uses}\ref{it:W11-3V},
\begin{equation}\label{eq:BY}
  B(Y)=\int_{\partial\Omega}Y\cdot\nu=3V\ne0 .
\end{equation}
For a given $X$ consider $E_{t,s}:=F^Y_s(F^X_t(\Omega))$.  By
Theorem~\ref{thm:first-var-volume} its volume is $C^1$ in $(t,s)$ with nonzero
$s$-derivative at the origin, so the implicit function theorem
\ref{base:calculus} gives a $C^1$ function $s=s(t)$ with $s(0)=0$ and
$|E_{t,s(t)}|=V$.  Differentiating minimality at $t=0$ and eliminating
$s'(0)$ gives $A(X)-\lambda B(X)=0$ with $\lambda:=A(Y)/B(Y)$; using
Theorems~\ref{thm:first-var-perimeter}, \ref{thm:first-var-volume} and
Lemma~\ref{lem:coulomb-first-var} to identify $A$ and $B$ gives
\eqref{eq:weak-EL}.
\end{proof}

\begin{remark}\label{rem:no-collar}
In a graph chart, vertical test fields reduce \eqref{eq:weak-EL} to the weak
prescribed-mean-curvature equation \eqref{eq:graph-PMC} below.  No smooth
collar is required at any point.
\end{remark}

\section{Regularity of the Newtonian potential}

\begin{lemma}[$C^{1,\beta}$ regularity of $v_\Omega$]\label{lem:potential-C1beta}
\uses{def:coulomb}
For bounded measurable $\Omega\subset\R^3$,
\begin{equation}\label{eq:potential-C1beta}
  v_\Omega\in C^{1,\beta}(\R^3)\qquad\text{for every }\beta<1 .
\end{equation}
\end{lemma}

\begin{proof}\uses{def:coulomb}
Differentiation under the integral is valid because $|x-y|^{-2}$ is locally
integrable:
\[
  \nabla v_\Omega(x)=-\int_\Omega\frac{x-y}{|x-y|^3}\,dy .
\]
Continuity follows by splitting the integral into a small ball and its
complement.  For the H\"older estimate let $h:=|x-z|<1/4$.  On
$B_{2h}(x)\cup B_{2h}(z)$ the integral of either gradient kernel is $O(h)$.
On the complement, the mean value theorem for $K(\xi):=\xi/|\xi|^3$ gives
$|K(x-y)-K(z-y)|\le Ch|x-y|^{-3}$, and since $\Omega$ is bounded, integrating
over dyadic annuli gives
\begin{equation}\label{eq:potential-log}
  |\nabla v_\Omega(x)-\nabla v_\Omega(z)|\le Ch(1+|\log h|).
\end{equation}
For every $\beta<1$, $h(1+|\log h|)\le C_\beta h^\beta$.
\end{proof}

\begin{fnote}[No Calder\'on--Zygmund]\label{fnote:no-CZ}
No second-derivative $L^p$ theorem is needed anywhere.
Lemma~\ref{lem:potential-C1beta} is proved directly, and
\S\ref{sec:green} uses tubular cutoffs rather than a Calder\'on--Zygmund
estimate.  Note also the logarithm in \eqref{eq:potential-log}: $v_\Omega$ is
\emph{not} claimed to be $C^{1,1}$, and no later step requires it.
\end{fnote}

\section{The smooth bootstrap}\label{sec:bootstrap}

\begin{lemma}[The graph equation]\label{lem:graph-PMC}
\uses{thm:constrained-first-var,not:minimizer-rep}
In a graph chart $x_3=f(x')$ over a planar disk, \eqref{eq:weak-EL} implies
the weak equation
\begin{equation}\label{eq:graph-PMC}
  \dvg\Bigl(\frac{\nabla f}{\sqrt{1+|\nabla f|^2}}\Bigr)
  =\pm\bigl(\lambda-v_\Omega(x',f(x'))\bigr),
\end{equation}
where the sign depends only on which side of the chart is $\Omega$.  Only
this direction is used later.
\end{lemma}

\begin{proof}\uses{thm:constrained-first-var,cor:graph-area}
Use a vertical field $X=\zeta(x')\eta(x_3)e_3$ with $\eta\equiv1$ near the
graph.  Writing $Q:=\sqrt{1+|\nabla f|^2}$, direct differentiation gives
$\dvgt X=(\nabla f\cdot\nabla\zeta)/Q^2$, while
$d\Hs^2=Q\,dx'$ and $X\cdot\nu\,d\Hs^2=\pm\zeta\,dx'$.  Substitution in
\eqref{eq:weak-EL} and planar integration by parts give
\eqref{eq:graph-PMC}.
\end{proof}

\begin{lemma}[Uniform ellipticity of the mean-curvature operator]
\label{lem:mc-elliptic}
For
\[
  a_{ij}(p):=\frac{\delta_{ij}}{\sqrt{1+|p|^2}}-\frac{p_ip_j}{(1+|p|^2)^{3/2}}
\]
and $|p|\le M$, the matrix $(a_{ij}(p))$ is uniformly elliptic with constants
depending only on $M$, and $p\mapsto a_{ij}(p)$ is smooth.
\end{lemma}

\begin{proof}
Diagonalise in the frame $\{p/|p|\}\cup p^\perp$: the eigenvalues are
$(1+|p|^2)^{-3/2}$ and $(1+|p|^2)^{-1/2}$.
\end{proof}

\begin{proposition}[$C^{2,\beta}$ from the quasilinear Schauder lemma]
\label{prop:bootstrap-C2}
\uses{lem:graph-PMC,thm:quasilinear-schauder,lem:potential-C1beta}
With $f$ as in Lemma~\ref{lem:graph-PMC} and $f\in C^{1,1/2}$, one has
$f\in C^{2,1/2}$; and then $f\in C^{2,\beta}$ for every $\beta<1$.
\end{proposition}

\begin{proof}
\uses{thm:quasilinear-schauder,lem:mc-elliptic,lem:potential-C1beta,
      thm:eps-regularity}
The $C^{1,1/2}$ regularity from Theorem~\ref{thm:eps-regularity} makes
$a_{ij}(\nabla f)$ uniformly elliptic with $C^{0,1/2}$ coefficients
(Lemma~\ref{lem:mc-elliptic}), and the right side of \eqref{eq:graph-PMC} is
$C^{0,1/2}$ by Lemma~\ref{lem:potential-C1beta}.  Apply
Theorem~\ref{thm:quasilinear-schauder} with $\alpha=1/2$ to get
$f\in C^{2,1/2}$.  For its full-norm estimate, $\|f\|_\infty$ is finite on
the fixed smaller graph chart; it is retained as an input, while the
derivative estimate is independent of an additive height constant.
Then $\nabla f$ is Lipschitz and
$v_\Omega\in C^{1,\beta}$, so the coefficients and the right side are
$C^{0,\beta}$ for every $\beta<1$; reapplying
Theorem~\ref{thm:quasilinear-schauder} gives $f\in C^{2,\beta}$.
\end{proof}

\begin{proposition}[The last derivative]\label{prop:bootstrap-C3}
\uses{prop:bootstrap-C2,thm:nondiv-schauder}
With $f$ as in Proposition~\ref{prop:bootstrap-C2}, $f\in C^{3,\beta}$ for
every $\beta<1$.  Consequently
\begin{equation}\label{eq:C3-boundary}
  \partial\Omega\in C^{3,\beta}\qquad\text{for every }\beta<1 .
\end{equation}
\end{proposition}

\begin{proof}\uses{prop:bootstrap-C2,thm:nondiv-schauder,lem:potential-C1beta}
One more derivative requires a separate line of argument: the divergence-form
estimate alone does not automatically add it.  Put
\[
  a_{ij}(x):=D_jA_i(\nabla f(x)),
  \qquad
  g(x):=\pm\bigl(\lambda-v_\Omega(x,f(x))\bigr),
\]
where $A_i(p):=p_i/\sqrt{1+|p|^2}$.  Then $a_{ij}\in C^{1,\beta}$ and
$g\in C^{1,\beta}$ by Propositions~\ref{prop:bootstrap-C2} and
Lemma~\ref{lem:potential-C1beta}, and for $w_k:=\partial_kf$,
differentiating the weak equation gives
\begin{equation}\label{eq:bootstrap-div}
  \partial_i\bigl(a_{ij}\partial_jw_k\bigr)=\partial_kg .
\end{equation}
Expanding the divergence --- legitimate now because $a\in C^{1,\beta}$ and
$w_k\in C^{1,\beta}$ --- yields the nondivergence equation
\begin{equation}\label{eq:bootstrap-nondiv}
  a_{ij}\partial_{ij}w_k=\partial_kg-(\partial_ia_{ij})\partial_jw_k ,
\end{equation}
whose right side is $C^{0,\beta}$.  Theorem~\ref{thm:nondiv-schauder},
including its harmless bounded first-order term, gives
$w_k\in C^{2,\beta}$.
\end{proof}

\begin{corollary}[Pointwise Euler--Lagrange equation]\label{cor:EL-pointwise}
\uses{prop:bootstrap-C3,conv:curvature}
On $\partial\Omega$,
\begin{equation}\label{eq:EL-pointwise}
  H+v_\Omega=\lambda .
\end{equation}
\end{corollary}

\begin{proof}\uses{prop:bootstrap-C3,lem:graph-PMC,conv:curvature}
By \eqref{eq:C3-boundary} the weak equation \eqref{eq:graph-PMC} may be
expanded pointwise, and the left side is $H$ in the sign convention of
Convention~\ref{conv:curvature}.
\end{proof}

\section{Connectedness and the scaling identity}

\begin{proposition}[Connectedness]\label{prop:connected}
\uses{not:minimizer-rep,prop:no-singular-points}
$\Omega$ is connected.
\end{proposition}

\begin{proof}
\uses{prop:no-singular-points,lem:locality,lem:coulomb-disjoint,
      lem:cross-decay,def:minimizer}
The compact smooth boundary has finitely many components
(Proposition~\ref{prop:no-singular-points}), so the bounded open set $\Omega$
has finitely many connected components.  If $\Omega$ were disconnected, write
$\Omega=A\mathbin{\dot\cup}B$ with $A,B$ nonempty unions of components.  Their
closures are disjoint: a common boundary point would have a graph neighbourhood
in which the interior side is connected, forcing the two components to
coincide.  Compactness then gives positive separation.

Translate $B$ by $te$.  For all large $t$, Lemma~\ref{lem:locality} gives
$\Per(A\cup(B+te))=\Per(A)+\Per(B)=\Per(\Omega)$, and the two
self-interactions are unchanged.  The original energy contains the strictly
positive cross-term $I(A,B)$ (Lemma~\ref{lem:coulomb-disjoint}), while
$I(A,B+te)\to0$ (Lemma~\ref{lem:cross-decay}).  A large translation is
therefore a same-volume competitor of strictly lower energy, contradicting
Definition~\ref{def:minimizer}.
\end{proof}

\begin{proposition}[Scaling identity for the multiplier]\label{prop:scaling-identity}
\uses{thm:constrained-first-var,lem:coulomb-scaling}
\begin{equation}\label{eq:scaling-identity}
  3V\lambda=2\Per(\Omega)+5D(\Omega)
  =5\Ec(\Omega)-3\Per(\Omega).
\end{equation}
\end{proposition}

\begin{proof}
\uses{thm:constrained-first-var,lem:coulomb-scaling,thm:first-var-perimeter,
      thm:first-var-volume,lem:coulomb-first-var}
Choose a compactly supported smooth field $X$ equal to $x$ on a neighbourhood
of the bounded set $\overline\Omega$; the compact cutoff is needed because $x$
itself is not compactly supported.  Its straight perturbation
$x\mapsto x+tX(x)$ agrees there exactly with dilation by $1+t$.  By
\eqref{eq:scaling},
$\Per(r\Omega)=r^2\Per(\Omega)$, $D(r\Omega)=r^5D(\Omega)$ and
$|r\Omega|=r^3V$, so stationarity of $\Ec-\lambda|\cdot|$ at $r=1$ ---
i.e.\ \eqref{eq:weak-EL} with this $X$ --- gives
$2\Per(\Omega)+5D(\Omega)=3V\lambda$.  The second form follows from
$\Ec=\Per+D$.
\end{proof}

\begin{definition}[Stationary domain]\label{def:stationary}
\uses{conv:curvature,def:coulomb}
A \emph{stationary domain of volume $V$ with multiplier $\lambda$} is a
bounded connected open $\Omega\subset\R^3$ with $|\Omega|=V$, boundary of
class $C^3$, satisfying \eqref{eq:EL-pointwise} pointwise on
$\partial\Omega$.
\end{definition}

\begin{corollary}[Minimisers are stationary]\label{cor:minimizer-stationary}
\uses{def:stationary,prop:scaling-identity}
Every fixed-volume minimiser at volume $V>0$ has a representative which is a
stationary domain of volume $V$, and it satisfies
\eqref{eq:scaling-identity}.
\end{corollary}

\begin{proof}
\uses{prop:no-singular-points,prop:bootstrap-C3,cor:EL-pointwise,
      prop:connected,prop:scaling-identity,lem:bounded-representative}
Collect Lemma~\ref{lem:bounded-representative},
Propositions~\ref{prop:no-singular-points}, \ref{prop:bootstrap-C3},
\ref{prop:connected}, Corollary~\ref{cor:EL-pointwise} and
Proposition~\ref{prop:scaling-identity}.
\end{proof}

